\section{Обновление существующего truncated-SVD solver:
low-rank correction $B=P+\Delta$}

\subsection{Цель}

Требуется модифицировать существующий solver в
\texttt{ADP/solver/SVD.py}.

Текущая реализация непосредственно строит малоранговую матрицу
\[
    B_r=A\Sigma V^\top,
    \qquad
    \rank(B_r)\le r<m,
\]
то есть ограничивает ранг самой матрицы $B$.

Необходимо заменить внутреннюю параметризацию на малоранговую
\emph{поправку} к текущему multi-index basis:
\[
\boxed{
    B=P+\Delta,
    \qquad
    \Delta=A\Sigma V^\top,
    \qquad
    \rank(\Delta)\le r.
}
\]

При этом
\[
    P\in\mathbb R^{m\times d},
    \qquad
    PP^\top=I_m,
\]
остаётся текущим $m$-мерным EDR basis.

Новый solver не должен уменьшать размерность EDR до $r$.
Параметр $r$ задаёт только максимальный ранг изменения текущего
подпространства за один HPAO solve.

Существующий внешний API
\[
    \texttt{solve(index\_init, U, I, mass, lambda\_prox, rank, ...)}
\]
желательно сохранить.

\subsection{Исходная задача}

При фиксированных локальных коэффициентах
\[
    g_j\in\mathbb R^m
\]
исходный HPAO functional имеет вид
\[
    F(B)
    =
    \sum_{j=1}^{J}
    N_j
    \left\|
        I_j-U_jB^\top g_j
    \right\|_2^2
    +
    \lambda
    \|B-P\|_F^2.
\]

Ввести
\[
    B=P+\Delta.
\]

Определить residual текущего basis:
\[
\boxed{
    e_j^{(0)}
    =
    I_j-U_jP^\top g_j.
}
\]

Тогда задача принимает вид
\[
\boxed{
    F(\Delta)
    =
    \sum_{j=1}^{J}
    N_j
    \left\|
        e_j^{(0)}
        -
        U_j\Delta^\top g_j
    \right\|_2^2
    +
    \lambda\|\Delta\|_F^2.
}
\]

Именно этот functional должен оптимизировать обновлённый solver.

\subsection{Внутреннее представление correction}

Хранить
\[
    \Delta
    =
    A\diag(s)V^\top
    =
    \sum_{k=1}^{K}
    s_k a_kv_k^\top,
\]
где
\[
    A\in\mathbb R^{m\times K},
    \qquad
    V\in\mathbb R^{d\times K}.
\]

В начале solve:
\[
\boxed{
    \Delta_0=0.
}
\]

Не использовать
\[
    B_0=0
\]
как текущую точку оптимизации.

Фактическая матрица HPAO на любом шаге:
\[
\boxed{
    B_k=P+\Delta_k.
}
\]

\subsection{Изменение инициализации}

В существующем внутреннем solver заменить смысл переменной
\texttt{B}.

Предпочтительный вариант:
\begin{itemize}
    \item явно переименовать её в \texttt{Delta};
    \item не хранить отдельно плотную матрицу $B=P+\Delta$;
    \item материализовывать её только при необходимости.
\end{itemize}

Инициализация должна иметь вид
\[
    A=\varnothing,
    \qquad
    V=\varnothing,
    \qquad
    s=\varnothing,
    \qquad
    \Delta=0.
\]

До rank-loop один раз вычислить
\[
\boxed{
    E^{(0)}
    =
    I-\operatorname{forward}(U,P,g).
}
\]

Здесь
\[
    E^{(0)}_j=e_j^{(0)}.
\]

\subsection{Residual внутри rank-loop}

Для текущей факторизации
\[
    \Delta
    =
    A\diag(s)V^\top
\]
correction prediction равен
\[
    U_j\Delta^\top g_j
    =
    \sum_k
    s_k
    (a_k^\top g_j)
    U_jv_k.
\]

Следовательно, residual необходимо вычислять как
\[
\boxed{
    e_j
    =
    e_j^{(0)}
    -
    U_j\Delta^\top g_j.
}
\]

При использовании существующего cache
\[
    W_{j,:,k}=U_jv_k
\]
получаем
\[
\boxed{
    e
    =
    E^{(0)}
    -
    \operatorname{einsum}
    \left(
        W,\,
        (gA)\odot s
    \right).
}
\]

Не использовать в этом месте исходный $I$ как residual относительно
нулевой матрицы.

\subsection{Rank-$1$ correction}

На очередном rank-шаге искать
\[
    \delta
    =
    \sigma av^\top,
    \qquad
    \|a\|_2=\|v\|_2=1.
\]

Обновление:
\[
    \Delta^+
    =
    \Delta+\sigma av^\top.
\]

Изменение objective должно быть
\[
\boxed{
    F(\Delta+\sigma av^\top)-F(\Delta)
    =
    D(a,v)\sigma^2
    -
    2R(a,v)\sigma.
}
\]

Здесь
\[
\boxed{
    D(a,v)
    =
    \sum_j
    N_j
    (a^\top g_j)^2
    \|U_jv\|_2^2
    +
    \lambda
}
\]
и
\[
\boxed{
    R(a,v)
    =
    \sum_j
    N_j
    (a^\top g_j)
    v^\top U_j^\top e_j
    -
    \lambda
    \langle
        \Delta,av^\top
    \rangle_F.
}
\]

Поэтому
\[
\boxed{
    \sigma^\star
    =
    \frac{R(a,v)}{D(a,v)}
}
\]
и predicted gain равен
\[
\boxed{
    \operatorname{gain}
    =
    \frac{R(a,v)^2}{D(a,v)}.
}
\]

Существующий критерий \texttt{rank\_tol} должен использовать этот gain.

\subsection{Изменение матрицы $Q$}

Для текущего residual $e_j$ определить
\[
    q_j
    =
    U_j^\top e_j.
\]

Data-gradient factor:
\[
    Q_{\mathrm{data}}
    =
    \sum_j
    N_jg_jq_j^\top.
\]

Так как
\[
    B-P=\Delta,
\]
регуляризационная часть gradient равна
\[
    \lambda\Delta.
\]

Следовательно, используемая в rank-$1$ поиске матрица должна быть
\[
\boxed{
    Q
    =
    Q_{\mathrm{data}}
    -
    \lambda\Delta.
}
\]

Эквивалентно, если код продолжает временно использовать
\[
    B=P+\Delta,
\]
можно вычислять
\[
    P-B=-\Delta.
\]

Необходимо проверить численно, что обе реализации совпадают.

\subsection{$a$-step}

При фиксированном $v$
\[
    R(a,v)=a^\top Qv.
\]

Определить
\[
\boxed{
    G_v
    =
    \sum_j
    N_j
    \|U_jv\|_2^2
    g_jg_j^\top
    +
    \lambda I_m.
}
\]

Тогда использовать существующий small solve:
\[
\boxed{
    a_{\mathrm{raw}}
    =
    G_v^{-1}Qv,
}
\]
\[
\boxed{
    a
    =
    \frac{a_{\mathrm{raw}}}
         {\|a_{\mathrm{raw}}\|_2}.
}
\]

Матрица имеет размер
\[
    m\times m.
\]

Существующую реализацию через Cholesky/\texttt{lstsq}
сохранить.

\subsection{$v$-step}

При фиксированном $a$ обозначить
\[
    \alpha_j=a^\top g_j.
\]

Необходимо решать
\[
\boxed{
\left[
    \sum_j
    N_j\alpha_j^2
    U_j^\top U_j
    +
    \lambda I_d
\right]
v_{\mathrm{raw}}
=
\sum_j
N_j\alpha_j
U_j^\top e_j
-
\lambda\Delta^\top a.
}
\]

Не формировать матрицы
\[
    U_j^\top U_j
\]
явно.

Сохранить существующие два backend:
\begin{enumerate}
    \item direct Cholesky для достаточно малых $d$;
    \item matrix-free LSMR для больших $d$.
\end{enumerate}

После solve:
\[
\boxed{
    v
    =
    \frac{v_{\mathrm{raw}}}
         {\|v_{\mathrm{raw}}\|_2}.
}
\]

Существующий normal-residual certificate сохранить.

\subsection{Упрощение ridge target в $v$-step}

В старой параметризации rank-$r$ матрицы $B$ proximal term
задавал target
\[
    P^\top a.
\]

В correction-параметризации regularizer равен
\[
    \lambda\|\Delta\|_F^2.
\]

Для нового rank-$1$ направления около текущей $\Delta$ ridge target
определяется только текущей correction.

Поэтому при решении задачи непосредственно по новой компоненте
необходимо использовать
\[
\boxed{
    z
    =
    -\Delta^\top a.
}
\]

Эквивалентная нормальная система:
\[
    (A^\top A+\lambda I)v_{\mathrm{raw}}
    =
    A^\top y
    -
    \lambda\Delta^\top a.
\]

Не использовать
\[
    z=(P-B)^\top a
\]
с $B$, обозначающим только correction.
Если переменная $B$ всё ещё обозначает полную матрицу $P+\Delta$,
то старое выражение
\[
    (P-B)^\top a
\]
остаётся корректным.

Необходимо выбрать одну семантику переменных и не смешивать их.

\subsection{Compact QR/SVD}

Существующую compact compression сохранить.

После получения новой компоненты
\[
    \sigma av^\top
\]
имеем
\[
    \Delta^+
    =
    \Delta+\sigma av^\top.
\]

Если
\[
    \Delta=LR^\top,
\]
то после добавления новой компоненты:
\[
    L^+
    =
    [L,\sigma a],
    \qquad
    R^+
    =
    [R,v].
\]

Выполнить thin QR
\[
    L^+=Q_LR_L,
    \qquad
    R^+=Q_RR_R.
\]

Затем
\[
    R_LR_R^\top
    =
    \widetilde A
    \Sigma
    \widetilde V^\top.
\]

Новые факторы correction:
\[
\boxed{
    A
    =
    Q_L\widetilde A,
    \qquad
    V
    =
    Q_R\widetilde V,
}
\]
с сохранением не более заданного rank budget.

Полный SVD матрицы $m\times d$ не выполнять.

\subsection{Joint re-fit singular values}

После compact SVD необходимо переоценить все singular scales
при фиксированных $A,V$.

Пусть
\[
    \Delta
    =
    \sum_{k=1}^{K}
    s_ka_kv_k^\top.
\]

Ввести
\[
    X_{j,:,k}
    =
    (a_k^\top g_j)U_jv_k.
\]

Тогда требуется решить
\[
\boxed{
    \min_s
    \sum_j
    N_j
    \left\|
        e_j^{(0)}
        -
        \sum_k s_kX_{j,:,k}
    \right\|_2^2
    +
    \lambda\|s\|_2^2.
}
\]

Нормальная система:
\[
\boxed{
    Ms=b,
}
\]
где
\[
    M_{k\ell}
    =
    \sum_j
    N_j
    X_{j,:,k}^\top X_{j,:,\ell}
    +
    \lambda\delta_{k\ell},
\]
\[
\boxed{
    b_k
    =
    \sum_j
    N_j
    X_{j,:,k}^\top e_j^{(0)}.
}
\]

В отличие от старой параметризации здесь
\textbf{не должно быть} дополнительного члена
\[
    \lambda a_k^\top Pv_k
\]
в правой части.

Причина:
\[
    \|B-P\|_F^2
    =
    \|\Delta\|_F^2,
\]
поэтому prior уже полностью учтён ridge-членом
\[
    \lambda I.
\]

\subsection{Новый factorized objective}

Заменить старую формулу objective для факторов на
\[
\boxed{
    F
    =
    \sum_j
    N_j
    \left\|
        e_j^{(0)}
        -
        \sum_k
        s_k
        (a_k^\top g_j)
        U_jv_k
    \right\|_2^2
    +
    \lambda
    \sum_k s_k^2.
}
\]

При ортонормированных $A,V$
\[
    \|\Delta\|_F^2
    =
    s^\top s.
\]

Новая функция должна иметь смысл
\[
\texttt{\_correction\_objective\_from\_factors(...)}.
\]

Удалить из неё старые члены
\[
    \|P\|_F^2
\]
и
\[
    -2\langle B,P\rangle_F.
\]

\subsection{Выход solver}

После окончания rank-recovery получить
\[
    \Delta_r
    =
    A\diag(s)V^\top.
\]

Сформировать
\[
\boxed{
    B_{\mathrm{raw}}
    =
    P+\Delta_r.
}
\]

Не использовать старую процедуру
\[
    \texttt{\_complete\_basis(V,s,P)}
\]
для correction solver.

Вместо неё ортонормировать полный обновлённый basis:
\[
    B_{\mathrm{raw}}^\top
    =
    QR,
\]
\[
\boxed{
    P_{\mathrm{new}}=Q^\top.
}
\]

После этого выполнить существующий local refit:
\[
    g_j^{\mathrm{new}}
    =
    \arg\min_g
    \left\|
        I_j-U_jP_{\mathrm{new}}^\top g
    \right\|_2^2.
\]

Вернуть
\[
    \texttt{HPAOResult}
    \left(
        P_{\mathrm{new}},
        g^{\mathrm{new}},
        \texttt{diagnostics}
    \right).
\]

\subsection{Сохранение старого solver}

Существующий вариант
\[
    \rank(B)\le r
\]
не удалять до окончания сравнительных экспериментов.

Необходимо поддержать два режима, например:
\[
\texttt{low\_rank\_target = "matrix"}
\]
для старого solver и
\[
\texttt{low\_rank\_target = "correction"}
\]
для нового.

Допустим также отдельный entry point
\[
    \texttt{solve\_correction(...)}.
\]

Предпочтителен вариант, при котором общие численные primitives
не дублируются:
\begin{itemize}
    \item \texttt{\_FlatU};
    \item direct $v$ solve;
    \item LSMR $v$ solve;
    \item normal-residual certificate;
    \item \texttt{\_solve\_small};
    \item compact QR/SVD.
\end{itemize}

\subsection{Diagnostics}

Добавить:
\[
    \texttt{low\_rank\_target="correction"},
\]
\[
    \texttt{correction\_rank},
\]
\[
    \texttt{correction\_singular\_values},
\]
\[
    \texttt{correction\_frobenius\_norm},
\]
\[
    \texttt{relative\_correction\_norm}
    =
    \frac{\|\Delta\|_F}{\|P\|_F},
\]
\[
    \texttt{rank\_gain\_history},
\]
\[
    \texttt{rank\_objective\_history}.
\]

Сохранить существующие:
\begin{itemize}
    \item LSMR iteration counters;
    \item normal-residual certificates;
    \item inner convergence;
    \item $U$-pass counter;
    \item direct/iterative solve counters.
\end{itemize}

\subsection{Обязательные invariant tests}

Добавить следующие тесты.

\paragraph{Zero correction.}

При
\[
    r=0
\]
или при остановке до принятия первой компоненты:
\[
    P_{\mathrm{new}}=P
\]
с точностью до basis orientation.

\paragraph{Penalty.}

Проверить численно:
\[
    \|B-P\|_F^2
    =
    \|\Delta\|_F^2.
\]

\paragraph{Factor objective.}

Для случайных малых задач сравнить factorized expression с прямым
вычислением
\[
    \sum_j
    N_j
    \|I_j-U_j(P+\Delta)^\top g_j\|^2
    +
    \lambda\|\Delta\|_F^2.
\]

Разность должна быть на уровне floating-point error.

\paragraph{Rank-1 scale.}

Для фиксированных $a,v$ проверить, что
\[
    \sigma^\star=R/D
\]
совпадает с минимумом одномерной задачи, полученным прямым solver.

\paragraph{Joint scale refit.}

После решения
\[
    Ms=b
\]
gradient objective по $s$ должен быть близок к нулю.

\paragraph{Monotonicity.}

Каждый принятый deterministic rank-step должен удовлетворять
\[
    F(\Delta_{k+1})
    \le
    F(\Delta_k)
    +
    \varepsilon_{\mathrm{num}}.
\]

\paragraph{Full-rank reference.}

На малых задачах сравнить correction solver с полным HPAO solver.
Для возрастающего $r$ ошибка objective должна в среднем уменьшаться.

Не требовать монотонности projector error по $r$ как математического
инварианта.

\subsection{Главный эксперимент}

Для каждого outer HPAO step сохранить полное reference-решение
\[
    B_\star
\]
и определить
\[
    \Delta_\star=B_\star-P.
\]

Посчитать singular values
\[
    \sigma_1(\Delta_\star)
    \ge
    \cdots
    \ge
    \sigma_m(\Delta_\star).
\]

Измерить
\[
\boxed{
    E_r
    =
    \frac{
        \sum_{k=1}^r
        \sigma_k(\Delta_\star)^2
    }{
        \sum_{k=1}^m
        \sigma_k(\Delta_\star)^2
    }.
}
\]

Сравнить:
\begin{enumerate}
    \item текущий full HPAO;
    \item существующий solver с $\rank(B)\le r$;
    \item новый solver с $\rank(B-P)\le r$.
\end{enumerate}

Для каждого варианта измерить:
\[
    F,
    \qquad
    \text{wall time},
    \qquad
    \text{$U$ passes},
    \qquad
    \text{effective rank},
\]
\[
    \text{projector distance to full HPAO},
    \qquad
    \text{projector distance to truth}
\]
при наличии ground truth.

\subsection{Следующий этап: stochastic extension}

Стохастическую версию не смешивать с первой реализацией correction solver.

Сначала получить deterministic solver
\[
    \rank(B-P)\le r
\]
при использовании всех $J$ центров.

После прохождения correctness tests добавить sampling
\[
    \mathcal B_t
    \subset
    \{1,\ldots,J\},
    \qquad
    |\mathcal B_t|=b.
\]

Для data-частей использовать множитель
\[
    \frac{J}{b}.
\]

Regularization
\[
    \lambda\|\Delta\|_F^2
\]
не масштабировать.

На stochastic rank-$1$ шаге использовать
\[
\widehat R_t(a,v)
=
\frac{J}{b}
\sum_{j\in\mathcal B_t}
N_j
(a^\top g_j)
v^\top U_j^\top e_j
-
\lambda
\langle\Delta,av^\top\rangle_F,
\]
\[
\widehat D_t(a,v)
=
\frac{J}{b}
\sum_{j\in\mathcal B_t}
N_j
(a^\top g_j)^2
\|U_jv\|^2
+
\lambda,
\]
\[
\boxed{
    \widehat\sigma_t
    =
    \frac{\widehat R_t}{\widehat D_t}.
}
\]

Обновление:
\[
    \Delta
    \leftarrow
    \Delta
    +
    \widehat\sigma_tav^\top.
\]

После превышения rank budget выполнить
\[
\boxed{
    \Delta\leftarrow\Pi_r(\Delta)
}
\]
через существующий compact QR/SVD.

Стохастический режим реализовывать только после подтверждения
эквивалентности deterministic correction objective прямому вычислению.
